% ============================================================================
% Chapter 25 --- Wave 7: TAC-aware audit and IR
% ----------------------------------------------------------------------------
% Purpose : enrich audit pipeline with TAC correlation; refresh IR runbooks
%           against the architecture's actual operational behaviour observed
%           in Wave 5 pilot.
% Page target: ~5 pages.
% Dependencies: Waves 2, 4, 5 required (D17.3); reference layer Ch. 12;
%               canonical schema D12.2; TAC-session sampling D12.8.
% Mirrors wave template (D18.2).
% ----------------------------------------------------------------------------
% Decisions registered in _coherence-EN.md:
%   D25.1 IR runbooks refreshed against the architecture's emergent behaviour
%   D25.2 SOC trained on TAC-correlated investigation queries
%   D25.3 enriched correlation is additive --- legacy paths preserved
%   D25.4 exit criteria reference Observability conformance subset reinforced
%         with audit-replay exercises (D16.5 Wave 7 mapping)
% ============================================================================

\chapter{Wave~7 --- TAC-aware audit and IR}
\label{ch:wave7-audit-ir}

\begin{caveatbox}[Dependencies]
\noindent Wave~7 applies the wave chapter template established in
\trchap{ch:wave0-foundations}{}. Required pre-requisites are
Wave~2 (observability baseline), Wave~4 (policy plane producing
decision logs) and Wave~5 (pilot demonstrating end-to-end
correlation). It applies the canonical TAC event schema of
\trchap{ch:observability}{} (\S\ref{sec:ob-schema}, D12.2) and
the TAC-session sampling rule (D12.8). Closure is gated by the
Observability conformance subset of \trchap{ch:conformance}{}
reinforced with audit-replay exercises.
\end{caveatbox}

Wave~7 transforms the observability pipeline from a passive
collection layer into an active investigative resource. The TAC
correlation that the canonical schema makes possible becomes the
primary axis of forensic and operational investigation; the
incident-response runbooks of the institution are refreshed against
the architecture's actual operational behaviour, observed during
the Wave~5 pilot. The wave is consequently both a technical
configuration effort and an organisational training exercise.

%-----------------------------------------------------------------------------
\section{Goal and dependencies}
\label{sec:w7-goal}

The goal of Wave~7 is to operationalise three capabilities that
the preceding waves have made technically possible but
organisationally unexercised:

\begin{itemize}[leftmargin=1.5em,nosep]
  \item \textbf{TAC-correlated investigation.} The Security
    Operations team (whether internal, contracted, or
    distributed) is trained to investigate incidents along the
    TAC correlation axis: starting from a suspect activity,
    retrieving the full activation chain across layers, and
    bounding the scope of the incident by TAC session rather than
    by host or by user.
  \item \textbf{Refreshed IR runbooks.} Existing
    incident-response runbooks are updated to incorporate the
    new investigative paths and the new containment options
    (TAC suspension, per-TAC re-evaluation, scoped revocation)
    that the architecture admits.
  \item \textbf{Audit replay.} The capability to replay
    consequential decisions against their recorded inputs and
    policy versions becomes part of routine audit rather than
    only of post-incident analysis.
\end{itemize}

%-----------------------------------------------------------------------------
\section{Pre-requisites}
\label{sec:w7-prereq}

Beyond Waves~2, 4 and 5, four operational conditions strengthen
Wave~7:

\begin{itemize}[leftmargin=1.5em,nosep]
  \item the Wave~5 pilot has produced a structured wave-closure
    report identifying the cross-layer behaviours that the
    runbooks must accommodate;
  \item the Security Operations function is engaged at the
    workshop level rather than at the consultation level: the
    team that will use the new runbooks has authored or co-
    authored them;
  \item the canonical TAC event schema is operating across the
    layers exercised by the pilot, with a verified
    \texttt{correlation\_id} propagation;
  \item the institution's existing incident-response framework
    (CSIRT or CSIRT-equivalent function, MoUs with
    \loc{national-gov-csirt} and the \gls{nren}-CSIRT where applicable)
    is documented
    so that the new runbooks integrate with it rather than
    supersede it.
\end{itemize}

%-----------------------------------------------------------------------------
\section{Coexistence pattern}
\label{sec:w7-coexistence}

Wave~7 is structurally additive: existing incident-response
runbooks remain authoritative throughout the wave, and the
TAC-aware extensions are added as enrichments rather than as
replacements. The pattern reduces the operational risk of the
wave to acceptable levels --- a wave that introduces new
investigative procedures cannot afford to disrupt the
institution's incident-response capacity.

\paragraph{Per-runbook update.} Each runbook is updated as a
self-contained exercise: the existing version is preserved, an
enriched version is authored, the two are exercised in parallel
on tabletop scenarios, and the enriched version replaces the
existing one only when the SOC team has demonstrated proficiency
on the new investigative path. The order of runbook update is
typically determined by the Wave~5 pilot's findings: the runbooks
most exercised by the pilot's scope are updated first.

\paragraph{Detection rules.} Detection rules in the analytical
store (whether SIEM-native or expressed in the open store of
\trchap{ch:observability}{}) are refactored to consume the
canonical schema. Existing rules are not deleted; they remain
parallel until the refactored rules are demonstrated equivalent
on a representative observation window.

%-----------------------------------------------------------------------------
\section{Trajectory: greenfield}
\label{sec:w7-greenfield}

For institutions in a greenfield position, Wave~7 is the wave in
which the IR function itself is built around the TAC correlation
model from the start. There is no incumbent runbook corpus to
update, but the absence is itself a challenge: the institution
must author a complete runbook set without the incremental
benefit of an existing baseline. Indicative duration is six to
nine months, dominated by the runbook authoring and tabletop
exercising rather than by technical configuration.

%-----------------------------------------------------------------------------
\section{Trajectory: brownfield single-suite-centric}
\label{sec:w7-brownfield-ms}

For institutions whose existing audit pipeline is built around
a cloud-vendor-native SIEM (for example Microsoft Sentinel), Wave~7
enriches the incumbent SIEM's investigative surface with the
canonical TAC schema while keeping it as the primary SOC tool.

\paragraph{Pattern.} The OpenTelemetry collector layer of Wave~2
is configured to emit the canonical TAC schema both into
the incumbent SIEM and into the parallel open store. SOC analysts
continue to operate the incumbent SIEM as their primary
investigative interface; the new investigative queries are authored
as the incumbent SIEM's detection content (e.g.\ Sentinel KQL)
consuming the canonical schema, with parallel queries authored
against the open store for portability. Audit replay --- the
capability to retrieve a recorded decision and its policy version
--- is exercised against the open store, since the policy plane's
decision log of Wave~4 was routed there without going through the
incumbent SIEM's per-GB ingestion premium.

\paragraph{What changes in Wave~7.} The incumbent SIEM's detection
corpus (e.g.\ Sentinel KQL) is refactored to use the canonical
schema fields rather than vendor-native field names; the runbooks
reference the canonical investigation paths; the SOC training
programme is refreshed.

\paragraph{What does not change.} The incumbent SIEM remains the
SOC's primary surface; the existing notification and escalation
channels with the institution's CSIRT and with
\loc{national-gov-csirt} counterparts are preserved.

\paragraph{Indicative duration.} Six to twelve months.

%-----------------------------------------------------------------------------
\section{Trajectory: brownfield hybrid}
\label{sec:w7-brownfield-hybrid}

For institutions whose existing audit pipeline is built around
Splunk or a comparable per-volume-priced analytical store,
Wave~7 produces a parallel pattern with one important
distinction: the volume management of Wave~2's collector layer
has by now demonstrated its operational discipline, and Wave~7
formalises the corresponding cost management.

\paragraph{Pattern.} The collector layer of Wave~2 routes
canonical events to Splunk for the SOC's primary investigative
surface, and to the open store for retention beyond the
cost-managed Splunk window. Detection content in Splunk is
refactored to consume canonical fields. Audit replay queries
the open store for the longer-horizon material that Splunk no
longer holds.

\paragraph{What changes in Wave~7.} The split between Splunk's
investigative window and the open store's longer retention
becomes operationally formalised --- runbooks for incidents
older than the Splunk window are explicitly authored against the
open store. The TAC-session sampling rule of D12.8 is applied
to bound the Splunk ingestion volume without losing correlation
within preserved sessions.

\paragraph{What does not change.} Splunk remains the SOC's
primary surface within its window; the institution's existing IR
relationships are preserved.

\paragraph{Indicative duration.} Six to twelve months.

%-----------------------------------------------------------------------------
\section{Reversibility criteria}
\label{sec:w7-reversibility}

Wave~7 is structurally additive and is reversible at the
runbook level throughout the wave. Beyond the wave, reversal is
operationally possible but is treated as a runbook re-authoring
exercise rather than as a wave rollback.

\begin{itemize}[leftmargin=1.5em,nosep]
  \item Per-runbook reversal: a runbook whose enriched version
    proves operationally inferior is reverted to the existing
    version, with the divergence diagnosed before any further
    runbook is updated.
  \item Detection content reversal: refactored detection rules
    that produce false-positive rates inconsistent with the
    SOC's tolerance are reverted to the prior version while the
    rule is re-tuned.
  \item The canonical schema itself is not subject to wave-level
    reversal (it was frozen by D12.2 at the architecture level).
\end{itemize}

%-----------------------------------------------------------------------------
\section{Exit criteria}
\label{sec:w7-exit}

Wave~7 closes when the Observability conformance subset of
\trchap{ch:conformance}{} (\S\ref{sec:cf-suite}) passes against
the enriched pipeline, when the audit-replay exercises succeed
on a representative sample, when the agreed runbook set has
been updated, and when the SOC has demonstrated proficiency on
the new investigative paths.

\begin{itemize}[leftmargin=1.5em,nosep]
  \item Observability conformance subset passes (canonical
    schema enforced, retention verified, cross-layer correlation
    succeeds, alternative-store re-routing succeeds);
  \item Audit replay: a sampled set of decisions from
    \trchap{ch:policy}{}'s log is replayed against recorded
    inputs and policy versions, with verdicts identical;
  \item Runbook coverage: the institutionally agreed set of
    runbooks has been updated; the remaining runbooks are in a
    documented schedule;
  \item SOC proficiency: tabletop exercises on TAC-correlated
    investigation have been conducted with documented outcomes,
    and at least one exercise has been conducted with
    institutional CSIRT or \gls{nren}-CSIRT participation.
\end{itemize}

\begin{realworldbox}[Audit and IR practice in research-intensive institutions]
\noindent Patterns reported in NREN community fora and at
REN-ISAC indicate that research-intensive institutions
typically operate IR functions in close coordination with
their \gls{nren}-CSIRT counterparts and with sector incident-sharing
networks. Specific runbook structures and exercise cadences
vary substantially across institutions and should be developed
locally rather than imported from external precedents. The
institutional CSIRT relationships, where they exist, would
benefit from being engaged early in the wave rather than at
its closure.
\end{realworldbox}

%-----------------------------------------------------------------------------
\section{Regulatory anchoring}
\label{sec:w7-regulatory}

\noindent Wave~7 is the wave at which the institution becomes
\emph{operationally accountable} for the audit-and-IR
obligations of the regulatory framework. It is the principal
anchor of \loc{cybersecurity-baseline-regime}
(in the EU, the national NIS2 transposition --- NIS2~Art.~23:
incident reporting, with a 24-hour early warning, 72-hour
notification and one-month final report) and of the
personal-data-breach notification duty of
\loc{data-protection-regime} (in the EU, \gls{gdpr}~Art.~33--34:
notification to \loc{national-dpa} and to data
subjects), with the
TAC-aware audit pipeline as the timeline-reconstruction
substrate. The wave also operationalises that regime's
incident-handling and effectiveness-assessment duties
(in the EU, NIS2~Art.~21(2)(b) and Art.~21(2)(f), the latter
through audit-replay exercises). It is the
operational anchor of 27001~A.8.15 (logging) and A.8.16
(monitoring activities) at maturity, and of NIST CSF
\textsf{DETECT} (DE) and \textsf{RESPOND} (RS) functions in
their entirety. Detailed mapping in
\trchap{ch:regulatory-mapping}{} \S\ref{sec:rm-gdpr},
\S\ref{sec:rm-nis2}, \S\ref{sec:rm-iso},
\S\ref{sec:rm-nist-csf}.
